\subsection{Equations of Functions} 
\begin{example}Consider the relation $f=\Set{(x,y)\mid y=x^2-5}$\lbr
\makeblankfill{We can tell that this relation is a function because }\lbr\makeblankfill{}\lbr
If we want to evaluate this function, what do we do?
\begin{itemize}
\item To find $f(2)=\makeblanksmall$, we find \makeblanksmall when \makeblank[1in]\makeblanknewf
\item To find $f(-3)=\makeblanksmall$, we find \makeblanksmall when \makeblank[1in]\makeblanknewf
\item To find $f(0)=\makeblanksmall$, we find \makeblanksmall when \makeblank[1in]\makeblanknewf
\end{itemize}
We notice a pattern, which allows us to make the following shorthand.
\end{example}
\begin{definition}[Function Notation]
The following notations describe the same relation.
$$f=\Set{(x,y)\mid y=\mbox{expression in $x$ only}}$$
$$f(x)=\mbox{expression in $x$}$$
A relation that can be expressed this way is always a function.
\end{definition}
\begin{example}
Let $g(x)=-5x+3$.
\begin{itemize}
\item Write this in relation notation: \makeblank
\item \makeblankfill{Evaluate $g(4)$: }
\end{itemize}
\end{example}
\begin{note}Earlier, we named relations with capital letters from the later part of the alphabet ($R$, $S$, $T$, etc). When writing function notation, we will typically label functions with lower-case letters from the beginning of the alphabet ($f$, $g$, $h$, $k$, etc).\end{note}